\begin{abstract}
Graph Neural Networks (GNNs) have shown promise in computer-aided diagnosis from multiple ultrasound images. However, their patient-level predictions do not reveal which images contribute to the result. We present an explanation framework that adapts Local Interpretable Model-agnostic Explanations (LIME) to graph-level prediction by perturbing each patient's image set. Instead of superpixels, we sample image subsets, reconstruct their graphs and estimate image-level influence. To generate informative perturbations, we propose two-stage Adaptive Class-Balanced Sampling: Stage I samples randomly, and Stage II favors the singleton-predicted image pool associated with the underrepresented subset class. Ridge and cross-validated Elastic Net quantify conditional influence, while Pearson correlation describes marginal association. We evaluated the framework in 135 positive-label patients with 3,072 ultrasound images from Taiwan Biobank. Adaptive sampling moved predicted class proportions closer to balance in 126 patients, but random sampling yielded lower mean Ridge probability-surrogate error (\RandomMAE{} versus \AdaptiveMAE{}). For both sampling schemes, deleting images in descending coefficient order lowered model output more than random deletion in all patients. The framework provides image-level visual explanations and a statistical basis for examining how sampling, approximation and ranking contribute to interpretation of multi-image predictions.
\end{abstract}
